Here $G$ is the derived cometric of the gauge field of Fact \claimref{sheaf.gauge-field-cometric}, $G(a,b) = \hb^{T}(a) \cdot \hb^{T}(b)$, as in Fact \claimref{synthesis.imported-gtg-dictionary}(i). All statements are finite exact computations from Fact \claimref{sheaf.gauge-field-cometric} and are verified as rational-function identities in $u$ by the script; we record the two on which the theorem depends. The Gram entries: $G(\rho,\rho) = r\,(\beta^2 - 1) = 2M - r = m^2 - r$; $G(\rho, \dtPG) = r^{1/2}(-\beta) = -m$; $G(\dtPG,\dtPG) = 1$; $G(\mu_i,\mu_i) = r^2 \cdot (-r^{-2}) = -1$; all cross terms with the angular sector vanish because the metric is block diagonal. The dying eigenvalue: the $(\rho, \dtPG)$ block has determinant $(m^2 - r)\cdot 1 - m^2 = -r$, and directly $g_k = (m^2 - r) + 2m(-m) + m^2 = -r$. That the \emph{eigenvalue}, not merely the entry, is first order follows from $\det G = -r$ together with the nonvanishing $3 \times 3$ minor being bounded away from zero at $E$: the dying eigenvalue is $\det G$ divided by a unit (the distinction matters; Example \ref{ex:eigenentry}). Note $g_k = -r$ holds \emph{independently of $M$}: the mass enters the radical's direction, not its rate. For (4): orthogonality $G(\dtS, \rho) = 0$ forces the stated transition coefficient, and substituting gives the stated $k$-decomposition whose $\rho$-part vanishes at $r = 0$. For (5): $k(\partial_t) = m$ and $k(\partial_r) = r^{1/2}$, so pairing with $\partial_t - \beta \partial_r$ gives $m - r^{1/2}\cdot m r^{-1/2} = 0$.
